\newpage
\scp{\tsc{Sula-Buru} classification}{09-03}
\tsc{Sula-Buru} is defined by:
*ə > \it{e}; *R > **h; *-ay/*-aw/*-a > \it{a}; *h/*q > {\0};
*-mp-/ *{\nbhyp}mb-/*ma-p/*ma-b > **mb.
Proto-Sula-\ili{Buru} has the following retentions,
some of which distinguish it from
neighbouring groups:
*ŋ = **ŋ; *k = **k; *d = **d; *y = **y; *z = **z; *b = **b; *p = **p

\vspace{5mm}
\dirtree{%%The percent after { is needed
%.1 {\tsc{Sula-Buru} *ə > \it{e}; *R > **h; *-ay/*-aw/*-a > \it{a}; *h/*q > {\0}; *-mp-/*-mb-/*ma-p/*ma-b > **mb (*ŋ = **ŋ; *k = **k; *d = **d; *y = **y; *z = **z; *b = **b; *p = **p)}.
.1 {\tsc{Sula-Buru}}.
		.2 {\tsc{Nuclear Sula-Buru} *z/*y > **y; *b > \it{f}; **mb > \it{b}; *d > **r; **muyu \tsc{neg}}.
				.3 {\tsc{Sula} *e > \it{i} /{\gap}{\#}; **h > \it{ʔ,} {\0}; *l > {\0} /Vα{\gap}Vα (**y = \it{y})}.
						.4 {\ili{Mangole} [mqc] **r > \it{l}; *uCu > \it{uCo}}.
						.4 {\ili{Sanana} [szn] **r > \it{h}; *uCu > \it{uCa}}.
				.3 {\tsc{Buru} **y > {\0}} (**r = \it{r}).
						.4 {\ili{Buru} [mhs] *uCu > \it{oCo, uCu}}.
						.4 {\ili{Lisela} [lcl] (*uCu = \it{uCu}; lexical and morphological influence from \ili{Kayeli} and \ili{Sanana})}.
		.2 {\ili{Ambelau} [amv] *p > \it{f}; *b > \it{b} /{\#}{\gap}, \it{β} /V{\gap}V; *w > \it{β}; **mb > \it{p}; **d/*l > \it{l, h}; *z > \it{l}; *ŋ/*n > \it{n}; *k > \it{k, ʔ,} {\0}; *t > \it{ʧ, r} /i{\gap}, \it{t} /n{\gap}, \it{r}; *y > traces}.
}
\hspace{3.5mm}--- \ili{Hukumina} [huw] (classification uncertain)